\documentclass[14pt]{extarticle}
% =============================================================================
% Shared Preamble for Worksheets, Homeworks, Exams & Lesson Notes
% (scripts/tex_preamble.tex)
% =============================================================================

\newif\ifdetailed
\ifdefined\setdetailed
  \setdetailed
\fi

\newif\ifwithsols
\ifdefined\setwithsols
  \setwithsols
\fi

\newif\ifwhiteblack
\ifdefined\setwhiteblack
  \setwhiteblack
\fi

% Document-class specific packages (loaded before common_preamble.tex / polyglossia)
\usepackage{fancyhdr}
\usepackage{enumitem}
\usepackage{titlesec}
\usepackage[bottom]{footmisc}
\usepackage{titling}
\usepackage{longtable} % page-breakable tables
\usepackage{float}
\usepackage[onehalfspacing]{setspace}

\ifwhiteblack
  \usepackage[a4paper, landscape, top=1.25cm, bottom=1.25cm, left=1.25cm, right=1.25cm, headheight=1.5cm, headsep=0.5cm]{geometry}
\else
  \usepackage[a4paper, top=2.4cm, bottom=2.4cm, left=2.2cm, right=2.2cm, headheight=1.4cm, headsep=0.4cm, footskip=1.0cm]{geometry}
\fi

% Shared common packages, fonts, polyglossia, callout boxes, code helpers, date & header macros
\IfFileExists{scripts/common_preamble.tex}{%
  % =============================================================================
% Shared Common Preamble (scripts/common_preamble.tex)
% Included by:
%   - scripts/tex_preamble.tex    (worksheets, homeworks, exams, lesson notes)
%   - scripts/book_preamble.tex   (two-sided Arabic book)
%   - scripts/beamer_preamble.tex (Beamer presentation slides)
% =============================================================================

\usepackage{fontspec}
\usepackage{amsmath}
\usepackage{amssymb}
\usepackage{xcolor}
\usepackage{graphicx}
\usepackage{listings}
\usepackage[most]{tcolorbox}
\usepackage{pifont}
\usepackage{minted}
\usepackage{etoolbox}
\usepackage{array}
\usepackage{booktabs}
\usepackage{tikz}
\usepackage{datetime2}

\usetikzlibrary{automata, positioning, arrows.meta, arrows, shapes.geometric, calc}

% Beamer already loads hyperref internally; load it here for article/book classes
% right before polyglossia (bidi).
\makeatletter
\@ifclassloaded{beamer}{}{%
  \usepackage{hyperref}%
  \hypersetup{hidelinks}%
  \providecommand{\LTcaptype}{table}%
}
\makeatother

% Load polyglossia (and bidi) AFTER graphics, tables, boxes, code, and hyperref
\usepackage{polyglossia}
\setmainlanguage[numerals=western]{arabic}
\setotherlanguage{english}
\newfontfamily\arabicfont[Script=Arabic]{Calibri}
\newfontfamily\arabicfontsf[Script=Arabic]{Calibri}
\newfontfamily\arabicfonttt[Script=Arabic]{Courier New}
\newfontfamily\emoji[BoldFont=*,ItalicFont=*]{Segoe UI Emoji}

% Allow slight inter-word stretch to avoid Overfull \hbox with inline \textenglish
\emergencystretch=1.5em

% -----------------------------------------------------------------------------
% Default Minted & Callout Box Definitions
% -----------------------------------------------------------------------------
\AtBeginEnvironment{minted}{\setstretch{0.95}}
\setminted{fontsize=\small, breaklines=true, vspace=2pt}

\tcbset{
  arc=2.5pt, outer arc=2.5pt,
  boxsep=1pt, left=6pt, right=6pt, top=3pt, bottom=3pt,
  toptitle=1.5pt, bottomtitle=1.5pt,
  before skip=6pt, after skip=6pt,
  fonttitle=\bfseries,
  fontupper=\setstretch{1.1}\setlength{\parskip}{2pt},
  colback=white, colframe=black, boxrule=0.6pt,
  boxcolor/.style 2 args={colback=#1, colframe=#2}
}

\newtcolorbox{boxDef}[1][]{boxcolor={blue!4!white}{blue!75!black},
  title={{\emoji📘} تعريف\ifx\\#1\\\else ~#1\fi :}}
\newtcolorbox{boxExercise}[1][]{boxcolor={cyan!4!white}{cyan!70!black},
  title={{\emoji🧩} تمرين\ifx\\#1\\\else ~#1\fi :}}
\newtcolorbox{boxExample}[1][]{boxcolor={yellow!6!white}{orange!85!black},
  title={{\emoji📝} مثال\ifx\\#1\\\else ~#1\fi :}}
\newtcolorbox{boxNote}[1][]{boxcolor={gray!7!white}{black!75},
  title={{\emoji✨} ملاحظة\ifx\\#1\\\else ~#1\fi :}}
\newtcolorbox{boxAttention}[1][]{boxcolor={magenta!5!white}{magenta!75!black},
  title={{\emoji🔔} تنبيه\ifx\\#1\\\else ~#1\fi :}}
\newtcolorbox{boxWarning}[1][]{boxcolor={red!4!white}{red!75!black},
  title={{\emoji⚡} ملاحظة هامة\ifx\\#1\\\else ~#1\fi :}}
\newtcolorbox{boxSolution}[1][]{boxcolor={green!4!white}{green!60!black},
  title={{\emoji✅} حل\ifx\\#1\\\else ~#1\fi :}}
\newtcolorbox{boxSymbol}[1][]{boxcolor={purple!4!white}{purple!70!black},
  title={{\emoji🔣} رمز\ifx\\#1\\\else ~#1\fi :}}
\newtcolorbox{boxHint}[1][]{boxcolor={teal!4!white}{teal!65!black},
  title={{\emoji💡} تلميح\ifx\\#1\\\else ~#1\fi :}}

% Compact code boxes: tighter line spacing and small vertical gaps
\tcbset{
  simplecode/.style={
    arc=2pt, outer arc=2pt,
    left=4pt, right=4pt, top=4pt, bottom=4pt,
    fontupper=\setstretch{0.95},
    before skip=6pt, after skip=6pt,
    colback=gray!5, colframe=black!50, boxrule=0.4pt
  }
}
\newenvironment{boxCode}{\begin{tcolorbox}[simplecode]}{\end{tcolorbox}}

% Code cell used inside tables (keeps code left-to-right inside RTL tables)
% Use inside a fixed-width p{...} column: \begin{pycell}\begin{minted}...\end{minted}\end{pycell}
\makeatletter
\newtcolorbox{pycell}{frame empty,colback=white,boxsep=0pt,left=2pt,right=2pt,top=2pt,bottom=2pt,width=\linewidth,before=\noindent,after=,baseline={\tcb@height-\ht\strutbox},fontupper=\setstretch{0.95}\setlength{\topsep}{0pt}\setlength{\partopsep}{0pt}\setlength{\parskip}{0pt}}
\makeatother

\newcolumntype{C}[1]{>{\centering\arraybackslash}p{#1}}

% Dotted date style (DD.MM.YYYY)
\makeatletter
\DTMnewdatestyle{dotted}{%
  \renewcommand{\DTMdisplaydate}[4]{\two@digits{##3}.\two@digits{##2}.##1}%
  \renewcommand{\DTMDisplaydate}{\DTMdisplaydate}%
}
\makeatother
\DTMsetdatestyle{dotted}

% -----------------------------------------------------------------------------
% Customizable Header / Footer Metadata
% -----------------------------------------------------------------------------
\newcommand{\headerleft}{مدرسة التسامح الشاملة}
\newcommand{\headerright}{الأستاذ محمود اغبارية}
\newcommand{\footercenter}{\thepage}

\newcommand{\setheaderleft}[1]{\renewcommand{\headerleft}{#1}}
\newcommand{\setheaderright}[1]{\renewcommand{\headerright}{#1}}
\newcommand{\setfootercenter}[1]{\renewcommand{\footercenter}{#1}}
%
}{%
  \IfFileExists{../../../scripts/common_preamble.tex}{%
    % =============================================================================
% Shared Common Preamble (scripts/common_preamble.tex)
% Included by:
%   - scripts/tex_preamble.tex    (worksheets, homeworks, exams, lesson notes)
%   - scripts/book_preamble.tex   (two-sided Arabic book)
%   - scripts/beamer_preamble.tex (Beamer presentation slides)
% =============================================================================

\usepackage{fontspec}
\usepackage{amsmath}
\usepackage{amssymb}
\usepackage{xcolor}
\usepackage{graphicx}
\usepackage{listings}
\usepackage[most]{tcolorbox}
\usepackage{pifont}
\usepackage{minted}
\usepackage{etoolbox}
\usepackage{array}
\usepackage{booktabs}
\usepackage{tikz}
\usepackage{datetime2}

\usetikzlibrary{automata, positioning, arrows.meta, arrows, shapes.geometric, calc}

% Beamer already loads hyperref internally; load it here for article/book classes
% right before polyglossia (bidi).
\makeatletter
\@ifclassloaded{beamer}{}{%
  \usepackage{hyperref}%
  \hypersetup{hidelinks}%
  \providecommand{\LTcaptype}{table}%
}
\makeatother

% Load polyglossia (and bidi) AFTER graphics, tables, boxes, code, and hyperref
\usepackage{polyglossia}
\setmainlanguage[numerals=western]{arabic}
\setotherlanguage{english}
\newfontfamily\arabicfont[Script=Arabic]{Calibri}
\newfontfamily\arabicfontsf[Script=Arabic]{Calibri}
\newfontfamily\arabicfonttt[Script=Arabic]{Courier New}
\newfontfamily\emoji[BoldFont=*,ItalicFont=*]{Segoe UI Emoji}

% Allow slight inter-word stretch to avoid Overfull \hbox with inline \textenglish
\emergencystretch=1.5em

% -----------------------------------------------------------------------------
% Default Minted & Callout Box Definitions
% -----------------------------------------------------------------------------
\AtBeginEnvironment{minted}{\setstretch{0.95}}
\setminted{fontsize=\small, breaklines=true, vspace=2pt}

\tcbset{
  arc=2.5pt, outer arc=2.5pt,
  boxsep=1pt, left=6pt, right=6pt, top=3pt, bottom=3pt,
  toptitle=1.5pt, bottomtitle=1.5pt,
  before skip=6pt, after skip=6pt,
  fonttitle=\bfseries,
  fontupper=\setstretch{1.1}\setlength{\parskip}{2pt},
  colback=white, colframe=black, boxrule=0.6pt,
  boxcolor/.style 2 args={colback=#1, colframe=#2}
}

\newtcolorbox{boxDef}[1][]{boxcolor={blue!4!white}{blue!75!black},
  title={{\emoji📘} تعريف\ifx\\#1\\\else ~#1\fi :}}
\newtcolorbox{boxExercise}[1][]{boxcolor={cyan!4!white}{cyan!70!black},
  title={{\emoji🧩} تمرين\ifx\\#1\\\else ~#1\fi :}}
\newtcolorbox{boxExample}[1][]{boxcolor={yellow!6!white}{orange!85!black},
  title={{\emoji📝} مثال\ifx\\#1\\\else ~#1\fi :}}
\newtcolorbox{boxNote}[1][]{boxcolor={gray!7!white}{black!75},
  title={{\emoji✨} ملاحظة\ifx\\#1\\\else ~#1\fi :}}
\newtcolorbox{boxAttention}[1][]{boxcolor={magenta!5!white}{magenta!75!black},
  title={{\emoji🔔} تنبيه\ifx\\#1\\\else ~#1\fi :}}
\newtcolorbox{boxWarning}[1][]{boxcolor={red!4!white}{red!75!black},
  title={{\emoji⚡} ملاحظة هامة\ifx\\#1\\\else ~#1\fi :}}
\newtcolorbox{boxSolution}[1][]{boxcolor={green!4!white}{green!60!black},
  title={{\emoji✅} حل\ifx\\#1\\\else ~#1\fi :}}
\newtcolorbox{boxSymbol}[1][]{boxcolor={purple!4!white}{purple!70!black},
  title={{\emoji🔣} رمز\ifx\\#1\\\else ~#1\fi :}}
\newtcolorbox{boxHint}[1][]{boxcolor={teal!4!white}{teal!65!black},
  title={{\emoji💡} تلميح\ifx\\#1\\\else ~#1\fi :}}

% Compact code boxes: tighter line spacing and small vertical gaps
\tcbset{
  simplecode/.style={
    arc=2pt, outer arc=2pt,
    left=4pt, right=4pt, top=4pt, bottom=4pt,
    fontupper=\setstretch{0.95},
    before skip=6pt, after skip=6pt,
    colback=gray!5, colframe=black!50, boxrule=0.4pt
  }
}
\newenvironment{boxCode}{\begin{tcolorbox}[simplecode]}{\end{tcolorbox}}

% Code cell used inside tables (keeps code left-to-right inside RTL tables)
% Use inside a fixed-width p{...} column: \begin{pycell}\begin{minted}...\end{minted}\end{pycell}
\makeatletter
\newtcolorbox{pycell}{frame empty,colback=white,boxsep=0pt,left=2pt,right=2pt,top=2pt,bottom=2pt,width=\linewidth,before=\noindent,after=,baseline={\tcb@height-\ht\strutbox},fontupper=\setstretch{0.95}\setlength{\topsep}{0pt}\setlength{\partopsep}{0pt}\setlength{\parskip}{0pt}}
\makeatother

\newcolumntype{C}[1]{>{\centering\arraybackslash}p{#1}}

% Dotted date style (DD.MM.YYYY)
\makeatletter
\DTMnewdatestyle{dotted}{%
  \renewcommand{\DTMdisplaydate}[4]{\two@digits{##3}.\two@digits{##2}.##1}%
  \renewcommand{\DTMDisplaydate}{\DTMdisplaydate}%
}
\makeatother
\DTMsetdatestyle{dotted}

% -----------------------------------------------------------------------------
% Customizable Header / Footer Metadata
% -----------------------------------------------------------------------------
\newcommand{\headerleft}{مدرسة التسامح الشاملة}
\newcommand{\headerright}{الأستاذ محمود اغبارية}
\newcommand{\footercenter}{\thepage}

\newcommand{\setheaderleft}[1]{\renewcommand{\headerleft}{#1}}
\newcommand{\setheaderright}[1]{\renewcommand{\headerright}{#1}}
\newcommand{\setfootercenter}[1]{\renewcommand{\footercenter}{#1}}
%
  }{%
    \input{../../scripts/common_preamble.tex}%
  }%
}

\pagestyle{fancy}

% Prevent widow/orphan lines
\widowpenalty=10000
\clubpenalty=10000

% Compact section headings and default list spacing so pages do not waste vertical space
\titlespacing*{\section}{0pt}{1.4ex plus 0.4ex minus 0.2ex}{0.7ex plus 0.2ex}
\titlespacing*{\subsection}{0pt}{1.1ex plus 0.3ex minus 0.2ex}{0.5ex plus 0.1ex}
\titlespacing*{\subsubsection}{0pt}{0.9ex plus 0.2ex minus 0.1ex}{0.4ex plus 0.1ex}
\setlist{topsep=4pt, partopsep=0pt, parsep=2pt, itemsep=4pt}

\gappto\captionsarabic{\renewcommand{\contentsname}{جدول المحتويات}}

% Legacy listings settings
\lstset{
  language=[Sharp]C,
  numbers=left,
  stepnumber=1,
  numbersep=8pt,
  frame=single,
  basicstyle=\ttfamily\small
}
\ifwhiteblack
  \lstset{
    keywordstyle=\bfseries,
    stringstyle=\ttfamily,
    commentstyle=\itshape
  }
\else
  \lstset{
    keywordstyle=\color{blue},
    stringstyle=\color{red},
    commentstyle=\color{green!50!black}
  }
\fi

% White-and-black landscape presentation mode overrides
\ifwhiteblack
  \AtBeginDocument{
    \LARGE
    \usemintedstyle{bw}
  }
  \setminted{vspace=2pt}
  \tcbset{
    colback=white, colframe=black, boxrule=1pt,
    boxcolor/.style 2 args={colback=white, colframe=black},
    simplecode/.append style={colback=white, colframe=black, boxrule=1pt}
  }
\fi

% Default to empty author to avoid "LaTeX Warning: No \author given" when omitted
\author{}

% Compact title block
\makeatletter
\renewcommand{\@maketitle}{%
  \begin{center}%
    {\ifwhiteblack\Huge\else\LARGE\fi \bfseries \@title \par}%
    \ifdefempty{\@author}{}{%
      \vskip 0.15em%
      {\large \@author \par}%
    }%
    {\small \hfill \textenglish{\DTMtoday} \par}%
  \end{center}%
  \vspace{-0.3cm}%
}
\makeatother

\apptocmd{\maketitle}{\thispagestyle{fancy}}{}{}

% Apply page headers and footer (using \headerleft, \headerright, \footercenter from common_preamble.tex)
\newcommand{\applyfancyheaders}{%
  \fancyhf{}%
  \fancyhead[L]{\textarabic{\headerleft}}%
  \fancyhead[R]{\textarabic{\headerright}}%
  \fancyfoot[C]{\textarabic{\footercenter}}%
}

\fancypagestyle{plain}{\applyfancyheaders}
\applyfancyheaders

\setcounter{tocdepth}{3} % only section, subsection, and subsubsection in TOC

\newcommand{\importpdfpage}[2]{%
    \newgeometry{left=0cm,right=0cm,top=0cm,bottom=0cm}
    \thispagestyle{empty}
    \includegraphics[page=#2,width=0.95\paperwidth,height=0.95\paperheight,keepaspectratio]{#1}
    \restoregeometry
}

\newcommand{\insertFullImg}[1]{%
    \noindent
    \makebox[\textwidth][c]{
        \includegraphics[width=0.9\paperwidth,keepaspectratio]{#1}%
    }%
}%


\ifwithsols
\title{حل وظيفة بيتية $8$ للصف العاشر: حلقات التكرار (\textenglish{while} و \textenglish{for})}
\else
\title{وظيفة بيتية $8$ للصف العاشر: حلقات التكرار (\textenglish{while} و \textenglish{for})}
\fi

\begin{document}

\maketitle
\thispagestyle{fancy}

\begin{enumerate}[itemsep=2em]

% --- من القسم الثالث: حلقات سهلة، مقدار الخطوة، ومعالجة مباشرة ---
\item اكتب برنامجًا يطبع جميع الأعداد الفردية من $99$ نزولًا إلى $33$ تنازليًّا، وبجانب كل عدد يطبع مضاعفه الثلاثي (أي العدد مضروبًا بـ $3$).
\ifdetailed
\begin{boxExample}
\begin{english}
\begin{minted}{text}
99 297
97 291
95 285
...
33 99
\end{minted}
\end{english}
\end{boxExample}
\fi
\ifwithsols
\begin{boxSolution}
\begin{english}
\begin{minted}{csharp}
for (int i = 99; i >= 33; i -= 2)
{
    Console.WriteLine(i + " " + (i * 3));
}
\end{minted}
\end{english}
\end{boxSolution}
\fi
\filbreak

\item اكتب برنامجًا يستقبل من المستخدم أسعار منتجات (أعدادًا حقيقية موجبة) في متجر خلال فترة تنزيلات، ويتوقف عند إدخال العدد $0$. \\
لكل سعر يُدخله المستخدم: إذا كان السعر أكبر من $200$ شيكل، يطبع البرنامج السعر بعد خصم $25\%$ منه (أي السعر مضروبًا بـ $0.75$)، وإلّا يطبع السعر كما هو بدون خصم.
\ifdetailed
\begin{boxExample}
\begin{english}
\begin{minted}{text}
Enter price (0 to stop):
120
Final price: 120
Enter price (0 to stop):
400
Final price: 300
Enter price (0 to stop):
0
\end{minted}
\end{english}
\end{boxExample}
\fi
\ifwithsols
\begin{boxSolution}
\begin{english}
\begin{minted}{csharp}
Console.WriteLine("Enter price (0 to stop):");
double price = double.Parse(Console.ReadLine());
while (price != 0)
{
    if (price > 200)
    {
        Console.WriteLine("Final price: " + (price * 0.75));
    }
    else
    {
        Console.WriteLine("Final price: " + price);
    }
    Console.WriteLine("Enter price (0 to stop):");
    price = double.Parse(Console.ReadLine());
}
\end{minted}
\end{english}
\end{boxSolution}
\fi
\filbreak

% --- من القسم الرابع: أسئلة المجموع ---
\item اكتب برنامجًا يستقبل من المستخدم عددًا صحيحًا موجبًا $n$ يمثّل عدد الأيام، ثم يستقبل عدد الساعات التي تدرّب فيها رياضي في كل يوم من الأيام الـ $n$ (أعدادًا صحيحة). \\
يحسب البرنامج ويطبع مجموع ساعات التدريب في \textbf{الأيام المكثّفة فقط} (اليوم المكثّف هو اليوم الذي تدرّب فيه الرياضي $3$ ساعات أو أكثر).
\ifdetailed
\begin{boxExample}
\begin{english}
\begin{minted}{text}
Enter number of days:
5
Enter training hours for 5 days:
2
4
1
3
5
Total intensive hours: 12
\end{minted}
\end{english}
\end{boxExample}
\fi
\ifwithsols
\begin{boxSolution}
\begin{english}
\begin{minted}{csharp}
Console.WriteLine("Enter number of days:");
int n = int.Parse(Console.ReadLine());
int sum = 0;

Console.WriteLine("Enter training hours for " + n + " days:");
for (int i = 1; i <= n; i++)
{
    int hours = int.Parse(Console.ReadLine());
    if (hours >= 3)
    {
        sum += hours;
    }
}
Console.WriteLine("Total intensive hours: " + sum);
\end{minted}
\end{english}
\end{boxSolution}
\fi
\filbreak

\item يريد طالب توفير مبلغ لشراء حاسوب محمول ثمنه $2500$ شيكل. اكتب برنامجًا يستقبل من الطالب المبالغ التي يوفّرها أسبوعيًّا (أعدادًا صحيحة موجبة) ويجمعها، ويتوقف فور وصول المجموع المتراكم إلى $2500$ شيكل أو تجاوزه. \\
بعد انتهاء الحلقة، يطبع البرنامج المبلغ الكلّي الذي تجمّع لدى الطالب، والمبلغ الفائض فوق ثمن الحاسوب.
\ifdetailed
\begin{boxExample}
\begin{english}
\begin{minted}{text}
Enter weekly savings:
800
700
600
550
Total saved: 2650
Extra amount: 150
\end{minted}
\end{english}
\end{boxExample}
\fi
\ifwithsols
\begin{boxSolution}
\begin{english}
\begin{minted}{csharp}
int sum = 0;
Console.WriteLine("Enter weekly savings:");
while (sum < 2500)
{
    int amount = int.Parse(Console.ReadLine());
    sum += amount;
}
Console.WriteLine("Total saved: " + sum);
Console.WriteLine("Extra amount: " + (sum - 2500));
\end{minted}
\end{english}
\end{boxSolution}
\fi
\filbreak

% --- من القسم الخامس: أسئلة العدّ ---
\item اكتب برنامجًا يستقبل من المستخدم $15$ عددًا صحيحًا ثلاثي المنزلة (بين $100$ و $999$)، ويحسب ويطبع كم عددًا منها رقم آحاده يساوي $5$ ويقبل القسمة على $3$ بدون باقٍ.
\ifdetailed
\begin{boxExample}
\begin{english}
\begin{minted}{text}
(For example, 105 and 225 are counted, while 115 and 123 are not)
\end{minted}
\end{english}
\end{boxExample}
\fi
\ifwithsols
\begin{boxSolution}
\begin{english}
\begin{minted}{csharp}
int count = 0;
for (int i = 1; i <= 15; i++)
{
    int num = int.Parse(Console.ReadLine());
    if (num % 10 == 5 && num % 3 == 0)
    {
        count++;
    }
}
Console.WriteLine("Count: " + count);
\end{minted}
\end{english}
\end{boxSolution}
\fi
\filbreak

\item اكتب برنامجًا يستقبل درجات حرارة يومية (أعدادًا صحيحة) ويتوقف عند إدخال القيمة $-100$ (قيمة التوقف لا تُحسب ضمن الأيام). \\
يحسب البرنامج ويطبع:
\begin{itemize}[nosep]
    \item عدد الأيام الكلّي التي أُدخلت درجات حرارتها.
    \item عدد الأيام الحارّة التي زادت فيها درجة الحرارة عن $30$ درجة.
\end{itemize}
\ifdetailed
\begin{boxExample}
\begin{english}
\begin{minted}{text}
Enter temperatures (-100 to stop):
25
33
29
36
31
-100
Total days: 5
Hot days: 3
\end{minted}
\end{english}
\end{boxExample}
\fi
\ifwithsols
\begin{boxSolution}
\begin{english}
\begin{minted}{csharp}
int totalDays = 0;
int hotDays = 0;

Console.WriteLine("Enter temperatures (-100 to stop):");
int temp = int.Parse(Console.ReadLine());
while (temp != -100)
{
    totalDays++;
    if (temp > 30)
    {
        hotDays++;
    }
    temp = int.Parse(Console.ReadLine());
}
Console.WriteLine("Total days: " + totalDays);
Console.WriteLine("Hot days: " + hotDays);
\end{minted}
\end{english}
\end{boxSolution}
\fi
\filbreak

% --- من القسم السادس: أسئلة القيمة القصوى والدنيا ---
\item اكتب برنامجًا يستقبل أزمنة $12$ متسابقًا في سباق جري (بالثواني، أعدادًا حقيقية موجبة)، ويجد ويطبع:
\begin{itemize}[nosep]
    \item زمن المتسابق الفائز بالمركز الأول (أي الزمن \textbf{الأصغر} من بين جميع المتسابقين).
    \item زمن المتسابق الأخير (أي الزمن \textbf{الأكبر} من بين جميع المتسابقين).
    \item الفرق الزمني بين المتسابق الأخير والفائز بالمركز الأول.
\end{itemize}
\ifwithsols
\begin{boxSolution}
\begin{english}
\begin{minted}{csharp}
double firstTime = double.Parse(Console.ReadLine());
double minTime = firstTime;
double maxTime = firstTime;

for (int i = 2; i <= 12; i++)
{
    double time = double.Parse(Console.ReadLine());
    if (time < minTime)
    {
        minTime = time;
    }
    if (time > maxTime)
    {
        maxTime = time;
    }
}
Console.WriteLine("Winner time (Min): " + minTime);
Console.WriteLine("Last place time (Max): " + maxTime);
Console.WriteLine("Difference: " + (maxTime - minTime));
\end{minted}
\end{english}
\end{boxSolution}
\fi
\filbreak

\item اكتب برنامجًا يستقبل أعدادًا صحيحة موجبة من المستخدم حتى إدخال العدد $-1$. \\
يجد البرنامج ويطبع \textbf{أكبر عدد يقبل القسمة على $10$} من بين الأعداد التي أُدخلت (افترض أنّ المستخدم سيدخل عددًا واحدًا على الأقل يقبل القسمة على $10$).
\ifdetailed
\begin{boxExample}
\begin{english}
\begin{minted}{text}
Enter positive numbers (-1 to stop):
45
120
95
80
150
63
-1
Max divisible by 10: 150
\end{minted}
\end{english}
\end{boxExample}
\fi
\ifwithsols
\begin{boxSolution}
\begin{english}
\begin{minted}{csharp}
int maxDiv10 = int.MinValue;
Console.WriteLine("Enter positive numbers (-1 to stop):");
int num = int.Parse(Console.ReadLine());

while (num != -1)
{
    if (num % 10 == 0 && num > maxDiv10)
    {
        maxDiv10 = num;
    }
    num = int.Parse(Console.ReadLine());
}
Console.WriteLine("Max divisible by 10: " + maxDiv10);
\end{minted}
\end{english}
\end{boxSolution}
\fi
\filbreak

% --- من القسم السابع: أسئلة متنوعة بمستوى متقدم قليلًا ---
% \item اكتب برنامجًا يستقبل من المستخدم وزن شخص على مدار $10$ أشهر متتالية (أعدادًا حقيقية موجبة)، ويحسب ويطبع في كم شهرًا \textbf{انخفض} وزن الشخص مقارنةً بالشهر الذي سبقه مباشرةً.
% \ifdetailed
% \begin{boxExample}
% \begin{english}
% \begin{minted}{text}
% (If the weights in the first 5 months are: 85, 83.5, 84, 82, 81
% then it decreased in month 2 (83.5 < 85), month 4 (82 < 84), and month 5 (81 < 82))
% \end{minted}
% \end{english}
% \end{boxExample}
% \fi
% \ifwithsols
% \begin{boxSolution}
% \begin{english}
% \begin{minted}{csharp}
% double prevWeight = double.Parse(Console.ReadLine());
% int decreaseCount = 0;

% for (int i = 2; i <= 10; i++)
% {
%     double currWeight = double.Parse(Console.ReadLine());
%     if (currWeight < prevWeight)
%     {
%         decreaseCount++;
%     }
%     prevWeight = currWeight;
% }
% Console.WriteLine("Months with weight decrease: " + decreaseCount);
% \end{minted}
% \end{english}
% \end{boxSolution}
% \fi
% \filbreak

% \item بطارية هاتف مشحونة بنسبة $100\%$. في كل ساعة يشغّل المستخدم تطبيقًا يستهلك نسبة معيّنة من البطارية. \\
% اكتب برنامجًا يستقبل في كل مرّة نسبة الاستهلاك بالساعة (عددًا صحيحًا موجبًا) ويطرحها من نسبة البطارية المتبقية. يستمرّ البرنامج طالما أنّ البطارية المتبقية أكبر من $15\%$ وطالما لم يتجاوز عدد ساعات الاستخدام $8$ ساعات. \\
% في النهاية يطبع البرنامج نسبة البطارية المتبقية وعدد ساعات الاستخدام الفعلية.
% \ifdetailed
% \begin{boxExample}
% \begin{english}
% \begin{minted}{text}
% Enter hourly battery usage:
% 20
% 25
% 30
% 15
% Remaining battery: 10%
% Hours used: 4
% \end{minted}
% \end{english}
% \end{boxExample}
% \fi
% \ifwithsols
% \begin{boxSolution}
% \begin{english}
% \begin{minted}{csharp}
% int battery = 100;
% int hours = 0;

% Console.WriteLine("Enter hourly battery usage:");
% while (battery > 15 && hours < 8)
% {
%     int usage = int.Parse(Console.ReadLine());
%     battery -= usage;
%     hours++;
% }
% Console.WriteLine("Remaining battery: " + battery + "%");
% Console.WriteLine("Hours used: " + hours);
% \end{minted}
% \end{english}
% \end{boxSolution}
% \fi
% \filbreak

\end{enumerate}

\vspace{1cm}
\begin{flushleft}
أرجو لكم وقتًا ممتعًا.

الأستاذ محمود اغبارية.
\end{flushleft}

\end{document}
